\chapter{Bias, varianza ed ensemble}
\label{cap:biasvar}

\begin{inbreve}
La SLT spiega underfitting e overfitting con la complessità del modello e il numero di dati. La decomposizione
bias-varianza dà un punto di vista complementare. Il training set è \emph{una} delle possibili realizzazioni dei
dati, e l'errore atteso si scompone in bias, varianza e rumore. Da qui si arriva agli ensemble (bagging,
boosting), che sfruttano la riduzione della varianza. Il capitolo chiude con qualche cenno sulla feature
selection.
\end{inbreve}

\section{L'idea}

Training set diversi, estratti dalla stessa distribuzione, producono ipotesi diverse. L'errore atteso in un
punto $\x$, mediato sui training set, si scompone in tre parti:
\begin{itemize}
  \item \textbf{bias}: la discrepanza tra la funzione vera e la predizione $h(\x)$ \emph{media}. È alto se $H$ è
        troppo piccolo, cioè se il modello è troppo rigido;
  \item \textbf{varianza}: la variabilità della risposta del modello al variare del training set. È dovuta a
        un'eccessiva flessibilità;
  \item \textbf{rumore}: le etichette contengono un errore casuale, quindi per uno stesso $\x$ ci possono essere
        più valori di $d$.
\end{itemize}
Si studia la regressione con target $d$ e loss quadratica.

\begin{esempio}[20 punti e 50 rette]
Sia $d = x + 2\sin(1{,}5x) + \eps$, con $\eps$ rumore gaussiano di media zero (varianza $0{,}2$). Si adatti ai
dati una retta $h(x) = w_1x+w_0$ ai minimi quadrati. Si osservano due fatti:
\begin{enumerate}
  \item la classe di ipotesi (rette) commette per questa $f$ un \textbf{errore sistematico};
  \item a seconda del dataset, i parametri trovati \textbf{sono diversi}. Ripetendo con 50 dataset di 20 punti si
        ottengono 50 rette diverse (\figrif{fig:bv-rette}).
\end{enumerate}
\end{esempio}

\begin{figure}[ht]
\centering
\begin{tikzpicture}
\begin{groupplot}[group style={group size=2 by 1, horizontal sep=0.5cm}, scatola, width=6.4cm, height=4.6cm,
  xmin=0, xmax=10, ymin=-2, ymax=13, xtick={0,5,10}, ytick={0,5,10}, xlabel={$x$}]
\nextgroupplot[ylabel={$d$}]
\addplot[thick, verde] table {dati/bv_vera.dat};
\addplot[only marks, mark=*, mark size=1.4pt, blu] table {dati/bv_punti.dat};
\addplot[thick, rosso] table {dati/bv_prima.dat};
\nextgroupplot[yticklabels={,,}]
\addplot[grigio!80, thin] table {dati/bv_rette.dat};
\addplot[thick, verde] table {dati/bv_vera.dat};
\addplot[very thick, rosso] table {dati/bv_media.dat};
\end{groupplot}
\end{tikzpicture}
\caption{A sinistra un dataset di 20 punti, la funzione vera (verde) e la retta adattata (rosso). A destra le 50
rette ottenute da 50 dataset diversi (grigio) e la loro media $\bar h$ (rosso). La distanza tra $\bar h$ e la
funzione vera è il bias, la dispersione delle rette attorno a $\bar h$ la varianza.}
\label{fig:bv-rette}
\end{figure}

\section{La decomposizione}
\label{sec:decomposizione}

Si assume che i dati siano estratti in modo \textbf{i.i.d.} da un'unica distribuzione $P$. Per un punto $\x$
arbitrario si vuole calcolare l'\textbf{errore di predizione atteso}
\[
\E_P\big[(d-h(\x))^2\big],
\]
dove $d$ è il target di $\x$ che potrebbe comparire in un dataset. Il simbolo $\E_P$ indica una media su due
fonti di casualità, entrambe governate da $P$:
\begin{itemize}
  \item il \textbf{training set}, estratto secondo $P$: a ogni training set corrisponde un'ipotesi $h$ diversa;
  \item il \textbf{rumore} nel target $d$ del punto $\x$, che è un dato nuovo, estratto indipendentemente dal
        training set.
\end{itemize}
Sia
\[
\bar h(\x) = \E_P[h(\x)]
\]
la \textbf{predizione media} in $\x$ delle ipotesi addestrate su dati estratti da $P$.

\begin{teorema}[Decomposizione bias-varianza]
\begin{equation}\label{eq:bias-varianza}
\boxed{\;\E_P\big[(d-h(\x))^2\big] = \underbrace{\E_P\big[(h(\x)-\bar h(\x))^2\big]}_{\text{varianza}}
+ \underbrace{\big(\bar h(\x)-f(\x)\big)^2}_{\text{bias}^2} + \underbrace{\E_P\big[(d-f(\x))^2\big]}_{\text{rumore}^2 \;=\;\sigma^2}\;}
\end{equation}
\end{teorema}
\begin{proof}
\begin{enumerate}
  \item \textbf{Si sviluppa il quadrato} e si usa la linearità del valore atteso:
        \[
        \E_P\big[(d-h(\x))^2\big] = \E_P\big[h(\x)^2 - 2d\,h(\x) + d^2\big] = \E_P\big[h(\x)^2\big] + \E_P\big[d^2\big] - 2\,\E_P[d\,h(\x)].
        \]
  \item \textbf{Termine misto.} $d$ è il target di un dato nuovo, mentre $h$ dipende solo dal training set. Il
        rumore del dato nuovo non dipende dal training set, quindi le due variabili sono indipendenti:
        $\E_P[d\,h(\x)] = \E_P[d]\,\E_P[h(\x)] = \E_P[d]\,\bar h(\x)$. Inoltre $\E_P[d] = \E_P[f(\x)+\eps] = f(\x)$,
        perché il rumore ha media zero. Il termine misto vale $-2f(\x)\bar h(\x)$.
  \item \textbf{Primo termine}, con il lemma della varianza (\rif{sec:probabilita}) applicato a $Z = h(\x)$:
        \[
        \E_P\big[h(\x)^2\big] = \E_P\big[(h(\x)-\bar h(\x))^2\big] + \bar h(\x)^2 .
        \]
  \item \textbf{Secondo termine}, con il lemma applicato a $Z = d$, la cui media è $f(\x)$:
        \[
        \E_P\big[d^2\big] = \E_P\big[(d-f(\x))^2\big] + f(\x)^2 .
        \]
  \item \textbf{Si raccoglie}:
        \begin{align*}
        \E_P\big[(d-h(\x))^2\big] &= \E_P\big[(h(\x)-\bar h(\x))^2\big] + \underbrace{\bar h(\x)^2 - 2f(\x)\bar h(\x) + f(\x)^2}_{=\,(\bar h(\x)-f(\x))^2} + \E_P\big[(d-f(\x))^2\big]\\
        &= \Var[h(\x)] + \Bias[h(\x)]^2 + \E_P[\eps^2] ,
        \end{align*}
        e $\E_P[\eps^2] = \sigma^2$ è la varianza del rumore. \qedhere
\end{enumerate}
\end{proof}

In parole: \textbf{errore di predizione atteso = varianza + bias$^2$ + rumore$^2$}.
\begin{itemize}
  \item \textbf{Bias} $(\bar h(\x)-f(\x))^2$: discrepanza tra la funzione vera e la predizione \emph{media}
        sui training set. È un \textbf{errore sistematico}, dovuto a uno spazio $H$ troppo piccolo o a un modello
        troppo rigido.
  \item \textbf{Varianza}: variabilità della risposta del modello sui diversi training set, dovuta a una
        \textbf{flessibilità eccessiva}.
  \item \textbf{Rumore}: anche la soluzione ottima sbaglia. È \textbf{irriducibile}: non dipende dal modello e
        rappresenta una tolleranza $\sigma$ sulla risposta.
\end{itemize}

\begin{figure}[ht]
\centering
\begin{tikzpicture}[scale=0.95]
\draw[thick, fill=grigiochiaro] plot[smooth cycle, tension=0.8] coordinates {(0,0) (3,-0.7) (6.6,-0.2) (7.2,2.0) (4.6,3.3) (1.2,3.0) (-0.4,1.6)};
\node[font=\small, grigio] at (5.6,-0.15) {spazio delle ipotesi $H$};
\draw[blu, thick, fill=blu!10] (2.0,1.2) ellipse (1.0 and 0.7);
\foreach \p in {(1.6,1.4),(2.4,1.0),(2.1,1.6),(1.5,0.9),(2.6,1.45),(1.9,0.75)} \fill[blu] \p circle (1.3pt);
\fill[inchiostro] (2.0,1.2) circle (1.8pt) node[below, font=\footnotesize, yshift=-9pt] {$\bar h$};
\fill[verde!70!black] (5.0,1.5) circle (2pt) node[below, font=\footnotesize] {ottimo in $H$};
\fill[rosso] (8.4,2.4) circle (2pt) node[right, font=\footnotesize] {$f$ (con rumore)};
\draw[<->, thick] (3.05,1.3) -- (4.92,1.5) node[midway, above, font=\footnotesize] {bias};
\draw[<->, thick, densely dashed] (5.08,1.55) -- (8.33,2.38) node[midway, above, font=\footnotesize, sloped] {rumore};
\draw[<->, blu] (1.0,0.1) -- (3.0,0.1) node[midway, below, font=\footnotesize] {varianza};
\end{tikzpicture}
\caption{Vista grafica della decomposizione. Training set diversi producono soluzioni diverse (punti blu, la cui
dispersione è la varianza). La loro media $\bar h$ dista dalla miglior soluzione ottenibile (bias). La soluzione
ottima stessa dista dai dati a causa del rumore.}
\label{fig:bv-grafica}
\end{figure}

\begin{esempio}[L'analogia delle freccette]
Il centro del bersaglio è la funzione vera e ogni freccetta un modello addestrato su un training set diverso
(\figrif{fig:freccette}). \textbf{Alto bias e bassa varianza} corrisponde all'\textbf{underfitting},
\textbf{basso bias e alta varianza} all'\textbf{overfitting}. L'analogia non è perfetta per il ML, ma aiuta.
\end{esempio}

\begin{figure}[ht]
\centering
\begin{tikzpicture}[scale=0.62]
\foreach \cx/\cy/\bx/\by/\s/\t in {0/0/0/0/0.25/{basso bias, bassa varianza}, 5/0/0/0/0.95/{basso bias, alta varianza}, 10/0/1.1/0.9/0.25/{alto bias, bassa varianza}, 15/0/1.0/0.8/0.95/{alto bias, alta varianza}} {
  \begin{scope}[xshift=\cx cm]
    \foreach \r in {2,1.4,0.8} \draw[grigio] (0,0) circle (\r);
    \fill[rosso!70] (0,0) circle (0.25);
    \pgfmathsetseed{7}
    \foreach \k in {1,...,9} \fill[blu] ({\bx+\s*rand},{\by+\s*rand}) circle (2.6pt);
    \node[font=\footnotesize, align=center, text width=3.2cm] at (0,-2.65) {\t};
  \end{scope}
}
\end{tikzpicture}
\caption{L'analogia delle freccette: il centro è la funzione vera, ogni punto un modello addestrato su un
training set diverso.}
\label{fig:freccette}
\end{figure}

\section{Bias, varianza e regolarizzazione}
\label{sec:bv-lambda}

Con la loss regolarizzata (\rif{sec:tikhonov})
\[
\mathrm{Loss}(\w) = \sum_p\big(d_p-h(\x_p)\big)^2 + \lambda\norm{\w}^2
\]
l'intercetta è tipicamente esclusa dalla penalità. Variando $\lambda$ si ottengono soluzioni meno regolarizzate
(complesse, $\lambda$ basso) o più regolarizzate (semplici, $\lambda$ alto). Nell'esperimento della
\figrif{fig:bv-lambda} si estraggono 25 dataset da 25 punti dalla sinusoide $\sin(2\pi x)$ con rumore. Su
ciascuno si addestra un modello con 24 funzioni di base gaussiane, per tre valori di $\lambda$.

\begin{figure}[ht]
\centering
\begin{tikzpicture}
\begin{groupplot}[group style={group size=2 by 3, horizontal sep=0.45cm, vertical sep=0.35cm}, scatola,
  width=6.2cm, height=3.6cm, xmin=0, xmax=1, ymin=-1.6, ymax=1.6, xtick={0,1}, ytick={-1,0,1},
  xticklabels={,}, yticklabels={,,}]
\nextgroupplot[ylabel={\footnotesize $\ln\lambda=2{,}6$}, title={\small 20 delle 25 ipotesi}]
\addplot[rosso!70, thin] table {dati/bv_lam0_fit.dat};
\nextgroupplot[title={\small media $\bar h$ e funzione vera}]
\addplot[thick, verde] table {dati/bv_sin.dat};
\addplot[thick, rosso] table {dati/bv_lam0_media.dat};
\nextgroupplot[ylabel={\footnotesize $\ln\lambda=-0{,}31$}]
\addplot[rosso!70, thin] table {dati/bv_lam1_fit.dat};
\nextgroupplot
\addplot[thick, verde] table {dati/bv_sin.dat};
\addplot[thick, rosso] table {dati/bv_lam1_media.dat};
\nextgroupplot[ylabel={\footnotesize $\ln\lambda=-2{,}4$}, xticklabels={$0$,$1$}]
\addplot[rosso!70, thin] table {dati/bv_lam2_fit.dat};
\nextgroupplot[xticklabels={$0$,$1$}]
\addplot[thick, verde] table {dati/bv_sin.dat};
\addplot[thick, rosso] table {dati/bv_lam2_media.dat};
\end{groupplot}
\end{tikzpicture}
\caption{Effetto di $\lambda$ su bias e varianza. A sinistra le ipotesi apprese da dataset diversi, a destra la
loro media (rosso) e la funzione vera (verde). Con $\lambda$ alto (in alto) le curve si somigliano, ma la media è
lontana dalla sinusoide. Con $\lambda$ basso (in basso) le curve sono molto diverse, ma la media coincide quasi
con la sinusoide.}
\label{fig:bv-lambda}
\end{figure}

\begin{itemize}
  \item \textbf{$\lambda$ alto}: \textbf{alto bias, bassa varianza}. La penalità domina, il modello dipende poco
        dai dati ed è troppo rigido.
  \item \textbf{$\lambda$ medio}: bias più basso, varianza più alta. Si dipende di più dal particolare training
        set.
  \item \textbf{$\lambda$ molto basso}: \textbf{basso bias, alta varianza}. L'errore di test di un singolo
        modello può essere alto (per la varianza), ma \emph{in media} le ipotesi sono corrette.
\end{itemize}
Un modello \textbf{sovra-regolarizzato} ha bias alto, uno \textbf{sotto-regolarizzato} varianza alta. Il minimo di
bias$^2$ + varianza coincide circa con il minimo dell'errore di test (\figrif{fig:bv-tradeoff}).

\begin{figure}[ht]
\centering
\begin{tikzpicture}
\begin{axis}[stile, width=8.4cm, height=5cm, xmin=-3, xmax=2.1, ymin=0, ymax=0.16, xlabel={$\ln\lambda$},
  xtick={-3,-2,-1,0,1,2}, ytick={0,0.05,0.1,0.15}, yticklabels={$0$,$0{,}05$,$0{,}10$,$0{,}15$}, legend pos=outer north east]
\addplot[thick, blu] table[x=lnl, y=bias2] {dati/bv_tradeoff.dat}; \addlegendentry{bias$^2$}
\addplot[thick, rosso] table[x=lnl, y=var] {dati/bv_tradeoff.dat}; \addlegendentry{varianza}
\addplot[thick, viola] table[x=lnl, y=somma] {dati/bv_tradeoff.dat}; \addlegendentry{bias$^2$ + varianza}
\addplot[thick, inchiostro, densely dashed, y filter/.expression={y-0.09}] table[x=lnl, y=test] {dati/bv_tradeoff.dat}; \addlegendentry{errore di test $-\ \sigma^2$}
\end{axis}
\end{tikzpicture}
\caption{Bias$^2$, varianza e loro somma al variare di $\ln\lambda$, stimati su 100 dataset. La curva tratteggiata
è l'errore di test medio meno la varianza del rumore ($\sigma^2 = 0{,}09$). Attenzione: qui la complessità è più
alta a \textbf{sinistra}, al contrario dei grafici usuali.}
\label{fig:bv-tradeoff}
\end{figure}

\begin{idea}[Il collegamento con la SLT]
Bias e varianza sono un'altra faccia del compromesso tra adattamento ai dati e complessità. Il bias corrisponde
all'incapacità del modello di rappresentare la funzione (alto $\Remp$, underfitting). La varianza corrisponde alla
sensibilità ai dati specifici (alta VC-confidence, overfitting). C'è una conseguenza pratica per la valutazione.
A parità di complessità l'errore di test varia molto da un training set all'altro, soprattutto per i modelli
complessi. Una singola misura può quindi essere ``fortunata'' o ``sfortunata'' (capitolo~\ref{cap:validazione}).
\end{idea}

\section{Ensemble learning}
\label{sec:ensemble}

L'idea degli \termine{ensemble} è combinare le risposte di \textbf{più modelli}. Nello schema più semplice (a
voto):
\begin{itemize}
  \item \textbf{regressione}: comitato a media semplice, $o(\x) = \frac1K\sum_{i=1}^Kh_i(\x)$;
  \item \textbf{classificazione}: voto tra molti classificatori (purché diano risposte \textbf{diversificate}),
        oppure si classifica dopo aver mediato le uscite continue.
\end{itemize}
In alternativa (\termine{stacking}) il combinatore è a sua volta un modello di ML.

\begin{teorema}[Il comitato non è peggiore della media dei suoi membri]
Per una loss convessa come l'errore quadratico, l'errore del comitato è minore o uguale alla media degli errori
dei singoli modelli:
\[
\mathrm{loss}\big(\E_i[h_i]\big) \le \E_i\big[\mathrm{loss}(h_i)\big]
\]
(è la disuguaglianza di Jensen).
\end{teorema}

\begin{esempio}[Due modelli]
Target $d=5$. Due modelli che predicono 2 e 4: media degli errori $\frac{(5-2)^2+(5-4)^2}{2} = 5$, errore del
comitato ($o=3$) pari a $(5-3)^2 = 4$. Con predizioni 6 e 4: media degli errori $\frac{1+1}{2}=1$, errore del
comitato $(5-5)^2 = 0$.
\end{esempio}

\subsection{Bagging}

Il \termine{bagging} (\ingl{bootstrap aggregating}) addestra $K$ modelli su \textbf{sottoinsiemi diversi} del
training set, ottenuti con il \textbf{bootstrap} (ricampionamento con reinserimento), e ne combina le risposte
(media in regressione, voto in classificazione). Il legame con la decomposizione è diretto (si ricordi il caso
``$\lambda$ basso''): modelli ad \textbf{alta varianza e basso bias} funzionano bene \emph{in media}.
\uline{La media delle $h$ riduce la varianza senza aumentare il bias.} Inoltre il bagging tende ad aumentare il
\textbf{margine} di un classificatore: la media di $K$ curve di separazione tende a stare ``nel mezzo''.

\subsection{Boosting}

Se i modelli commettono gli \textbf{stessi errori}, combinarli non serve. Il \termine{boosting} (ad esempio
AdaBoost) li differenzia di proposito:
\begin{enumerate}
  \item addestra i classificatori \textbf{in sequenza}, concentrandosi sugli \textbf{errori}: il secondo dà più
        peso alle istanze sbagliate dal primo (ad esempio sono più probabili nel suo training set), e così via;
  \item combina i risultati con un \textbf{voto pesato}, dando più peso ai classificatori con errore più basso.
\end{enumerate}
Il boosting usa \termine{weak learner}, cioè classificatori appena migliori del caso (con errore $<1/2$ nel caso
binario). Se non lo si ferma, impara a classificare correttamente \textbf{tutte} le istanze di training e
costruisce in modo incrementale modelli complessi. Può migliorare molto learner deboli. Riduce il bias senza
aumentare la varianza: per questo resiste all'overfitting anche quando l'errore di training arriva a zero. Come il
bagging, tende a produrre separazioni a margine ampio. La resistenza all'overfitting però vale solo se i dati sono
poco rumorosi. Il boosting \textbf{soffre con dati rumorosi}: gli esempi rumorosi sono sbagliati di continuo,
ricevono pesi sempre maggiori e il modello finisce per adattarsi al rumore. Per questo non risolve di per sé il
problema underfitting/overfitting.

\section{Feature selection (cenni)}
\label{sec:feature-selection}

Trovare il sottoinsieme delle feature più informative porta grandi benefici. C'è la \textbf{riduzione della
dimensionalità}, con effetti sulla VC-dimension e sulla maledizione della dimensionalità. C'è il filtraggio di
informazione irrilevante e di rumore. C'è l'\textbf{interpretabilità}. Però:
\begin{itemize}
  \item è \textbf{computazionalmente difficile}: i sottoinsiemi possibili sono moltissimi e ogni valutazione
        richiede di riaddestrare il modello;
  \item si usa tipicamente una \textbf{ricerca euristica} (greedy, algoritmi genetici o altre tecniche), che può
        includere o meno il learner;
  \item i risultati sono molto sensibili al metodo di validazione: attenzione al \textbf{feature subset
        selection bias}, cioè all'errore di fare la selezione sull'intero dataset (\rif{sec:controesempio}).
\end{itemize}

\begin{puntichiave}
\begin{itemize}
  \item $\E_P[(d-h(\x))^2] = \Var[h(\x)] + \Bias[h(\x)]^2 + \sigma^2$; il rumore è irriducibile.
  \item Modelli rigidi (o molto regolarizzati): alto bias, bassa varianza; modelli flessibili (poco regolarizzati):
        basso bias, alta varianza.
  \item Un comitato non è peggiore della media dei suoi membri. Bagging: media di modelli ad alta varianza, che
        riduce la varianza. Boosting: learner deboli in sequenza, concentrati sugli errori.
\end{itemize}
\end{puntichiave}

\begin{domande}
\begin{enumerate}
  \item Derivare la decomposizione bias-varianza per la loss quadratica.
  \item Cosa rappresentano bias, varianza e rumore? Quale componente è irriducibile?
  \item Come variano bias e varianza con il parametro di regolarizzazione $\lambda$?
  \item Perché un comitato non è peggiore della media dei suoi membri?
  \item Differenza tra bagging e boosting; perché il bagging riduce la varianza?
\end{enumerate}
\end{domande}
